\section{Methodology}

\subsection{Controlled Experimental Design}

Our design follows directly from the generalization void hypothesis: fix the base model, training dataset, evaluation harness, and random seeds --- vary only whether training uses SFT, RLEF-Fraction, or RLEF-Binary. This allows the difficulty-scaling question to be answered without confounds from model choice, data selection, or evaluation protocol.

\textbf{Research questions:}
\begin{itemize}
  \item \textbf{RQ1 (Difficulty scaling):} Does RLEF produce a statistically significant positive-ordered advantage over SFT across five benchmark difficulty levels?
  \item \textbf{RQ2 (Signal void):} Is SFT's failure at hard difficulty a dataset void or a generalization void?
  \item \textbf{RQ3 (Reward formulation):} Does fraction-of-tests reward produce meaningfully better performance than binary reward?
\end{itemize}

\textbf{Controlled variables:}

\begin{table}[h]
\centering
\caption{Controlled experimental variables.}
\label{tab:controlled}
\begin{tabular}{ll}
\toprule
\textbf{Variable} & \textbf{Fixed Value} \\
\midrule
Base model & DeepSeek-Coder-7B-base \cite{Guo2024DeepSeekCoder} \\
Training data & APPS \cite{Hendrycks2021APPS} \\
Evaluation harness & bigcode-evaluation-harness (correctness-only) \\
Optimizer & GRPO \cite{Shao2024DeepSeekMath}, custom SimpleGRPOTrainer \\
Random seed & 42 throughout \\
Hardware & NVIDIA H100 NVL (5$\times$ 95830 MiB) \\
\bottomrule
\end{tabular}
\end{table}

\subsection{Base Model and Training Data}

We use \textbf{DeepSeek-Coder-7B-base} for all experiments. APPS provides training problems across three difficulty levels: introductory ($\sim$3500), interview ($\sim$5000), and competition (361 with $\geq$1 test case). Only problems with executable test cases are included in RLEF training. A key property: APPS competition problems have \textbf{85.32\% reference solution coverage} (308/361) --- directly relevant to RQ2 (Section~\ref{sec:results}).

\subsection{Training Methods}

\textbf{SFT Baseline:} HuggingFace Trainer with causal language modeling loss on APPS reference solutions. Key hyperparameters: lr=$2\times10^{-5}$, 3 epochs (1 for smoke scale), effective batch size 16, max\_length=2048.

\textbf{RLEF-Fraction:} GRPO starting from the SFT checkpoint, reward = fraction of test cases passing. Implemented via a custom SimpleGRPOTrainer compatible with PyTorch 2.5.1 / TRL 1.x. Subprocess-sandboxed execution (3.0s timeout). Key hyperparameters: lr=$1\times10^{-6}$, G=4 per prompt, max\_new\_tokens=512 (critical --- see Section~\ref{sec:results}), $\beta$=0.04 (KL penalty), temperature=0.8.

\textbf{RLEF-Binary (Proxy):} Full independent RLEF-Binary training was not completed due to compute constraints. Binary performance is estimated from RLEF-Fraction output via per-problem conversion. This is an acknowledged limitation (Section~\ref{sec:discussion}).

\textbf{SimpleGRPOTrainer} computes group-relative advantages $A_i = (r_i - \text{mean}(r)) / (\text{std}(r) + \varepsilon)$ and updates with REINFORCE-style loss plus KL penalty from the SFT reference. RewardMonitorCallback logs per-step rewards by difficulty.

\subsection{Evaluation Benchmarks}

\begin{table}[h]
\centering
\caption{Evaluation benchmarks and difficulty levels.}
\label{tab:benchmarks}
\begin{tabular}{lll}
\toprule
\textbf{Benchmark} & \textbf{Problems} & \textbf{Difficulty} \\
\midrule
HumanEval \cite{Chen2021HumanEval} & 164 & Easy \\
MBPP \cite{Austin2021MBPP} & 374 & Medium-Easy \\
LiveCodeBench-Easy \cite{Jain2024LiveCodeBench} & $\sim$200 & Medium \\
LiveCodeBench-Medium & $\sim$200 & Medium-Hard \\
LiveCodeBench-Hard & $\sim$100 & Hard (primary target) \\
\bottomrule
\end{tabular}
\end{table}

All evaluation uses bigcode-evaluation-harness in correctness-only mode (pass@1, single generation, greedy decoding). Identical harness parameters across all checkpoints.

\subsection{Statistical Analysis}

\textbf{Jonckheere-Terpstra (JT) test:} Tests the ordered alternative hypothesis that $\Delta$ values are non-decreasing across five difficulty levels. Appropriate because it tests an \textit{a priori} ordered hypothesis without requiring strict monotonicity at every adjacent pair. Implemented via pairwise Mann-Whitney U statistic sums, bootstrapped (n=5000).

\textbf{Bootstrap confidence intervals:} 95\% BCa intervals for $\Delta$ and $\Delta_{\text{ratio}}$ estimates (n=5000).

\textbf{Caveat:} Statistical values are conditional on smoke-scale proxy estimates, not independent experimental replications. See Section~\ref{sec:discussion}, Limitation L4.

\subsection{Sub-Hypothesis Structure}

\begin{table}[h]
\centering
\caption{Sub-hypothesis structure and outcomes.}
\label{tab:hypotheses}
\begin{tabular}{llll}
\toprule
\textbf{Hypothesis} & \textbf{Focus} & \textbf{Gate} & \textbf{Result} \\
\midrule
h-e1 & RLEF mechanism validity & MUST\_WORK & PASS \\
h-m1 & SFT failure characterization & MUST\_WORK & PASS \\
h-m2 & Non-zero reward monitoring & SHOULD\_WORK & FAIL (artifact) \\
h-m3 & Fraction vs binary ablation & SHOULD\_WORK & FAIL (null result) \\
h-m4 & Difficulty-ordered trend & SHOULD\_WORK & PASS \\
\bottomrule
\end{tabular}
\end{table}
